\subsection{Implications for scaling software}
\paragraph{Separate numerical representation from computational benefit}
The fixed-basis comparisons show large changes in exported conditioning, much
of which subsequent equilibration removes; the kernel ranking can even reverse
(\cref{subsec:results-bridge}). File-level diagnostics therefore complement,
rather than predict, solver outcomes. The mechanism audit also requires a
kernel-matched comparison: exact GM local normalization removes the block-scale
information in the extrema statistic, leaving a Euclidean coupling kernel
(\cref{prop:audit-collapse}). On the locally whitened triangular model the optimal
common factor is $(1+\rho^2)^{-1/2}$ (\cref{prop:audit-triangular-optimum}).
These statements explain spectral behavior without attributing it to metadata
that no longer affects the selected factors.

The final-policy ablation gives a separate implementation-level result: on its
six-model subset the guards suppress coupling changes, making local-only and
full exports identical (\cref{subsec:results-ablation}). Thus a coupling kernel
explains one conditioning comparison, whereas local transformations change the
dispatch representation. Neither observation identifies a solver's internal
cause of improved effort.

\paragraph{Compare contracts with strong numerical alternatives}
Budgets express application-unit accuracy and coefficient compatibility,
including when factors stay unchanged. Such an audit establishes precision
interpretation, not a budget-induced speedup. The activation follow-up supplies
the complementary case where the contract changes factors
(\cref{subsec:results-activation}). Its useful precision--effort regime is
configuration-dependent; earlier settings retain bound contradictions or incur
substantial solver work. The direct controls in the results include global
tolerance adjustment, NumericFocus and the unscaled model. In exact arithmetic,
fixed factors also admit a sufficient common raw-residual requirement
$\epsilon\le\min_r\alpha_r\delta_r$; setting a nominal solver tolerance does not
prove this hypothesis. Contracts make the row-specific requirements explicit
and allow separate checks of export and integer-rounding effects.

\paragraph{Verify the complete pipeline}
Presolve, internal scaling, tolerances and finite precision intervene between
the written row and the returned point. Their combined effect varies across
the three tested solvers and can change after a one-ulp perturbation. Exact
feasible witnesses of the actual imported exports can refute reported bounds;
the absence of a detected contradiction cannot certify them. The secondary
floating-point screen in Online Resource~1, Table~S10 remains exploratory.
Integer-aware GCD reduction and pattern repair answer different questions from
equivalent row scaling and are evaluated as such. Developers should measure
precision, candidate quality and complete-process cost together, as the
following deployment guidance makes explicit.
Online Resource~1 also derives an explicit finite-export residual reserve:
finite variable bounds make it uniform; otherwise only a pointwise bound is
asserted. It is an elementary accounting identity, with separate rational checks,
rather than an additional scaling algorithm.
